\let\negmedspace\undefined
\let\negthickspace\undefined
\documentclass[journal]{IEEEtran}
\usepackage[a5paper, margin=10mm, onecolumn]{geometry}
%\usepackage{lmodern} % Uncomment if needed for pdflatex


\setlength{\headheight}{1cm} % Set the height of the header box
\setlength{\headsep}{0mm}     % Set the distance between the header box and the top of the text

\usepackage{gvv-book}
\usepackage{gvv}
\usepackage{cite}
\usepackage{amsmath,amssymb,amsfonts,amsthm}
\usepackage{algorithmic}
\usepackage{graphicx}
\usepackage{textcomp}
\usepackage{xcolor}
\usepackage{txfonts}
\usepackage{listings}
\usepackage{enumitem}
\usepackage{mathtools}
\usepackage{gensymb}
\usepackage{comment}
\usepackage[breaklinks=true]{hyperref}
\usepackage{tkz-euclide} 
\usepackage{listings}
%\usepackage{gvv}                                        
\def\inputGnumericTable{}                                 
\usepackage[latin1]{inputenc}                                
\usepackage{color}                                            
\usepackage{array}                                            
\usepackage{longtable}                                       
\usepackage{calc}                                             
\usepackage{multirow}                                         
\usepackage{hhline}                                           
\usepackage{ifthen}                                           
\usepackage{lscape}
\usepackage{tikz}
\usepackage{circuitikz}
\usepackage{standalone} % For including external TikZ files

\begin{document}

\bibliographystyle{IEEEtran}
\vspace{3cm}

\title{6.5.1.3}
\author{EE24BTECH11013 - MANIKANTA D}
% \maketitle
% \newpage
% \bigskip
{\let\newpage\relax\maketitle}

\renewcommand{\thefigure}{\theenumi}
\renewcommand{\thetable}{\theenumi}
\setlength{\intextsep}{10pt} % Space between text and floats

\numberwithin{equation}{enumi}
\numberwithin{figure}{enumi}
\renewcommand{\thetable}{\theenumi}
\textbf{Question}:\\
Find the Maximum and Minimum values of the function if exists \\
 \begin{align*}
     y (x) = -\brak{x-1}^2+10
 \end{align*}  
\textbf{Solution : }\\
\textbf{Theoretical Solution :}\\

Finding the derivative of the function\\
\begin{align}
    \frac{dy}{dx} = -2(x-1)
\end{align}
Setting $\frac{dy}{dx} = 0$, we get:
\begin{align}
    x = 1
\end{align}
At $x = 1$, substituting into $y(x)$:
\begin{align}
    y(1) = -\brak{1-1}^2 + 10 = 10
\end{align}
To determine the nature of this critical point, we find the second derivative:
\begin{align}
    \frac{d^2y}{dx^2} = -2
\end{align}
Since $\frac{d^2y}{dx^2} < 0$, the function has a maximum at $x = 1$.

The function is a downward-opening parabola, so there is no absolute minimum. The function decreases without bound as $x \to \pm\infty$.\\

\textbf{Computational Solution :}\\

Maximum value of the function can be computed using the \textbf{Gradient ascent} method:\\
\textbf{Update the position iteratively}:
\begin{align}
    x_{n+1} = x_{n} + \eta \cdot \frac{dy}{dx}
\end{align}
Taking $x_0 = 0$ as the starting point\\
 Here, $\eta = 0.1$, where $\eta$is the learning rate.
\begin{align}
    x_{n+1} = x_{n} + \eta \cdot \brak{-2(x_n-1)}
\end{align}
The gradient ascent method converges to $x \approx 1$, and the maximum value is:
\begin{align}
    y(1) \approx 10
\end{align}

Minimum value of the function cannot be computed because the function decreases indefinitely as $x \to \pm\infty$.\\

\textbf{Computational Results :}\\
- Absolute Maximum:
\begin{align}
    x \approx 1 , \text{ } y(x) \approx 10
\end{align}
\begin{figure}[ht!]
   \centering
   \includegraphics[width=\columnwidth]{Figure_1.png}
\end{figure}
\end{document}
